% !TeX program = xelatex

% compile with XeLaTeX
% this template was created by salim bou 
\documentclass[dvipsnames,mathserif, aspectratio=169]{beamer}
\usepackage{setspace, float}
\setstretch{1.2}
\usepackage{tikz}
\usepackage{subcaption, tasks}

\usepackage{amsmath, amssymb, amsthm, nicematrix, pst-node, nccmath, cancel}

\AtBeginDocument{
	\allowdisplaybreaks
	\renewcommand{\jot}{10pt}
	\abovedisplayskip=5pt
	\belowdisplayskip=5pt
	\renewcommand{\theequation}{\arabic{equation}}
}

\usepackage{polyglossia, quran}
\setdefaultlanguage[numerals=maghrib]{arabic} % locale=mashriq, libya, algeria, tunisia, morocco, or mauritania  for names of months in \date 
\setotherlanguage{english}
\newfontfamily\arabicfont[Script=Arabic]{Amiri}
\newfontfamily\arabicfontsf[Script=Arabic]{Amiri}
\newfontfamily{\timesfont}{Times New Roman}

\newcommand{\ar}{\textarabic}
\newcommand{\en}{\textenglish}

\usepackage[T1]{fontenc}
\usepackage{times}

\usepackage[lite, zswash]{mtpro2}

\DeclareMathSymbol{0}{\mathalpha}{operators}{`0}
\DeclareMathSymbol{1}{\mathalpha}{operators}{`1}
\DeclareMathSymbol{2}{\mathalpha}{operators}{`2}
\DeclareMathSymbol{3}{\mathalpha}{operators}{`3}
\DeclareMathSymbol{4}{\mathalpha}{operators}{`4}
\DeclareMathSymbol{5}{\mathalpha}{operators}{`5}
\DeclareMathSymbol{6}{\mathalpha}{operators}{`6}
\DeclareMathSymbol{7}{\mathalpha}{operators}{`7}
\DeclareMathSymbol{8}{\mathalpha}{operators}{`8}
\DeclareMathSymbol{9}{\mathalpha}{operators}{`9}

\newcommand{\N}{\mathbb{N}}
\newcommand{\Z}{\mathbb{Z}}
\newcommand{\Q}{\mathbb{Q}}
\newcommand{\R}{\mathbb{R}}
\newcommand{\C}{\mathbb{C}}
\newcommand{\LL}{\mathcal{L}}

\usetheme{CambridgeUS}
\setbeamertemplate{bibliography item}{\insertbiblabel}
%\usecolortheme{crane}

% for RTL liste
\makeatletter
\newcommand{\RTListe}{\raggedleft\rightskip\leftm}
\newcommand{\leftm}{\@totalleftmargin}
\makeatother



% RTL frame title
\setbeamertemplate{frametitle}
{\vspace*{-1mm}
	\nointerlineskip
	\begin{beamercolorbox}[sep=0.3cm,ht=2.2em,wd=\paperwidth]{frametitle}
		\vbox{}\vskip-2ex%
		\strut\hskip1ex\insertframetitle\strut
		\vskip-0.8ex%
	\end{beamercolorbox}
}


% align subsection in toc
\makeatletter
\setbeamertemplate{subsection in toc}
{\leavevmode\rightskip=5ex%
	\llap{\raise0.1ex\beamer@usesphere{subsection number projected}{bigsphere}\kern1ex}%
	\inserttocsubsection\par%
}
\makeatother

% RTL triangle for itemize
\setbeamertemplate{itemize item}{\scriptsize\raise1.25pt\hbox{\donotcoloroutermaths$\blacktriangleleft$}} 

%\setbeamertemplate{itemize item}{\rule{4pt}{4pt}}

\defbeamertemplate{enumerate item}{square2}
{\LR{
		%
		\hbox{%
			\usebeamerfont*{item projected}%
			\usebeamercolor[bg]{item projected}%
			\vrule width2.25ex height1.85ex depth.4ex%
			\hskip-2.25ex%
			\hbox to2.25ex{%
				\hfil%
				{\color{fg}\insertenumlabel}%
				\hfil}%
		}%
}}

\setbeamertemplate{enumerate item}[square2]

\setbeamertemplate{navigation symbols}{}

\let\definition\relax
\let\enddefinition\relax

\newtheorem{definition}{تعريف}

\let\example\relax
\let\endexample\relax
\newtheorem{example}{مثال}





\begin{document}
	\author{\textbf{الطالبة : رملة قاسم محمد}}
	\title{\textbf{طريقة التحويل التفاضلي لحل نموذج النمو السكاني}}
	\date{\textbf{اشراف\\
			 د. خالد عبدالاله}}
	
		
		\begin{frame}
			\maketitle
		\end{frame}
		
		\begin{frame}{الخلاصة}
\textbf{الفصل الأول:} في هذا الفصل عرّفنا بعض المفاهيم الأساسية في المعادلات التفاضلية الإعتيادية منها تعريف المعادلة التفاضلية ورتبة المعادلة والدرجة وكذلك طريقة رانج كتا من الرتبة الرابعة العددية لحل المعادلات التفاضلية.\\[10pt]
\textbf{الفصل الثاني:} في هذا الفصل عرّفنا التحويل التفاضلي وكيفية تطبيقه على المعادلات التفاضلية وإستخدمناه في حل نموذج النمو السكاني وكذلك قارنا طريقة التحويل التفاضلي مع طريقة رانج كتا من الرتبة الرابعة
		\end{frame}
		
		\begin{frame}
			\begin{center}
				\Huge
				\textbf{الفصل الأول}\\
				\textbf{مفاهيم أساسية}
			\end{center}
		\end{frame}
		
		\begin{frame}
			\begin{definition}[المعادلة التفاضلية]
				    هي علاقة بين المتغير المعتمد والمتغير المستقل (المتغيرات المستقلة) تدخل فيها المشتقات أو التفاضلات وتسمى المعادلة التفاضلية اعتيادية إذا كان المتغير المعتمد دالة في متغير مستقل واحد وبالتالي لا تحتوي إلا على مشتقات عادية. حيث الشكل العام لها
				\[
				F(x,y',y'',y''',\dots,y^{(n)})=0
				\]
				حيث $y^{(n)}=\mfrac{d^ny}{dx^n}$ المشتقة من الرتبة $n$ للمتغير $y$ بالنسبة إلى $x$.
			\end{definition}
		\end{frame}
		
		\begin{frame}{طريقة رانج كتا من الرتبة الرابعة}			
			تُستخدم طريقة رانج كتا من الرتبة الرابعة لحل المعادلات التفاضلية العادية من الشكل:
			\[
			\frac{dy}{dx} = f(x, y), \quad y(x_0) = y_0
			\]
			نفترض أننا نريد إيجاد الحل على الفترة \( [x_0, x_n] \)، فنقوم بتقسيمها إلى \( n \) أجزاء متساوية:
			\[
			h = \frac{x_n - x_0}{n}
			\]
			ونحصل على نقاط التقسيم:
			\[
			x_1 = x_0 + h,\quad x_2 = x_0 + 2h,\quad \dots,\quad x_n = x_0 + nh
			\]
			وتهدف طريقة رانج كتا لإيجاد قيم $y$ عند قيم $x$، وهي كالتالي:
			\begin{itemize}
				\item \( y_0 \): القيمة الابتدائية المعطاة عند \( x_0 \).
				\item \( y_1 \): تقريب لقيمة الحل عند النقطة \( x_1 \).
				\item \( y_2 \): تقريب لقيمة الحل عند \( x_2 \).
			\end{itemize}
		\end{frame}
		
		\begin{frame}
			
			أي أن \( y_k \) يمثل تقريبًا للقيمة الحقيقية للدالة \( y(x_k) \)، ويتم حساب هذه القيم بشكل تتابعي باستخدام الطريقة.
			لإيجاد \( y_{n+1} \) من \( y_n \) باستخدام خطوة \( h \)، نحسب:
			\[
			k_1 = f(x_n, y_n)
			\]
			\[
			k_2 = f\left(x_n + \frac{h}{2}, \, y_n + \frac{h}{2} k_1 \right)
			\]
			\[
			k_3 = f\left(x_n + \frac{h}{2}, \, y_n + \frac{h}{2} k_2 \right)
			\]
			\[
			k_4 = f(x_n + h, \, y_n + h k_3)
			\]
			ثم نحسب القيمة الجديدة:
			\[
			y_{n+1} = y_n + \frac{h}{6} \left( k_1 + 2k_2 + 2k_3 + k_4 \right)
			\]
		\end{frame}
		
		\begin{frame}
			\begin{center}
				\Huge
				\textbf{الفصل الثاني}\\
				\textbf{طريقة التحويل التفاضلي}
			\end{center}
		\end{frame}
		
		\begin{frame}{مقدمة}\doublespacing
			تعد طريقة التحويل التفاضلي (DTM) وسيلة عددية وتحليلية لحل المعادلات التكاملية  والمعادلات التفاضلية العادية والجزئية وأنظمة المعادلات التفاضلية، حيث توفر هذه الطريقة الحل على شكل سلاسل متقاربة ذات مكونات سهلة الحساب. وبالنسبة للمعادلات التفاضلية، تقوم هذه الطريقة ببناء حل تحليلي في صورة كثير حدود، وهي تختلف عن طريقة سلسلة تايلور التقليدية ذات الرتب العالية التي تتطلب حسابات رمزية للمشتقات الضرورية لدوال البيانات، إذ تستغرق طريقة سلسلة تايلور وقتاً طويلاً في الحسابات عند التعامل مع الرتب الكبيرة، بينما تمثل DTM إجراءً تكرارياً للحصول على حلول تحليلية لسلاسل تايلور للمعادلات التفاضلية
		\end{frame}
		
		\begin{frame}{وصف الطريقة}
			من  المعروف انه اذا كانت لدينا دالة مثل $u$ قابلة للاشتقاق بإستمرارية الى ما لا نهاية، فأنه يمكن كتابة الدالة $u$ على شكل متسلسلة تايلر كالتالي
			\begin{equation}
				u(x) = \sum_{k=0}^{\infty} \frac{1}{k!} \frac{d^k u(x_0)}{dx^k} (x - x_0)^k
			\end{equation}
			نعرف التحويل التفاضلي (Differential transform) للدالة $u$ من الرتبة $k$ ، والذي يرمز له بالرمز $U(k)$ ، كالتالي
			\begin{equation}
				U(k) = \frac{1}{k!} \frac{d^k u(x_0)}{dx^k}
			\end{equation}
			بينما معكوس التحويل التفاضلي يكون
			\begin{equation}
				u(x) = \sum_{k=0}^{\infty} U(k) (x - x_0)^k
			\end{equation}
		\end{frame}
		
		\begin{frame}
				\textbf{المثال الاول} هو محاولة لوصف النمو السكاني هي النموذج الذي قدمه Thomas Malthus عام 1798 ، والذي استند فيه إلى فرضية مفادها أن معدل نمو السكان يتناسب طردياً مع إجمالي عدد السكان $p(t)$ عند الزمن $t$، أي أن:
				\begin{equation}
					\left\{
					\begin{array}{l}
						\dfrac{dp}{dt} = kp\\
						p(0) = p_0
					\end{array}
					\right.
				\end{equation}
				حيث $k$ هو عامل التناسب. الحل التحليلي لهذه المعادلة التفاضلية:
				\begin{equation}
					p(t) = p_0 e^{kt}
				\end{equation}
				نوجد حل المعادلة (5) مع الشرط الابتدائي
				$p(0)=2$ و $k=1.1$ على الفترة الزمنية $[0,5]$ بإستعمال طريقة التحويل التفاضلي ومقارنته بالحل التحليلي
				$p(t) = 2 e^{1.1t}$
				\begin{equation}
					\left\{
					\begin{array}{l}
						\dfrac{dp}{dt} = 1.1p\\
						p(0) = 2
					\end{array}
					\right.
				\end{equation}
		\end{frame}
		
		\begin{frame}
			
			بتطبيق تحويل DTM نحصل على:
			\[
			(k+1)P(k+1) = 1.1\,P(k)
			\]
			ومنها:
			\[
			P(k+1) = \frac{1.1}{k+1} P(k)
			\]
			\textbf{من الشرط الابتدائي:}
			\[
			P(0)=2
			\]
			\textbf{إيجاد التكرارات:}
			\[
			P(1) = \frac{1.1}{1}P(0) = 1.1 \times 2 = 2.2
			\]
			\[
			P(2) = \frac{1.1}{2}P(1) = \frac{1.1}{2} \times 2.2 = 1.21
			\]	
			\[
			P(3) = \frac{1.1}{3}P(2) = \frac{1.1}{3} \times 1.21 = 0.4437
			\]
		\end{frame}
		
		\begin{frame}
			\[
			P(4) = \frac{1.1}{4}P(3) = \frac{1.1}{4} \times 0.4437 = 0.1220175
			\]
			\[
			P(5) = \frac{1.1}{5}P(4) = \frac{1.1}{5} \times 0.1220175 = 0.02684385
			\]
			\textbf{الحل التقريبي حتى $5$ حدود:}
			\[
			p(x) \approx \sum_{k=0}^{5} P(k)x^k
			\]
			\[
			p(x) \approx 2 + 2.2x + 1.21x^2 + 0.4437x^3 + 0.1220175x^4 + 0.02684385x^5
			\]
		\end{frame}
		
		\begin{frame}{الحل بطريقة \en{RK4}}
			
			نريد حل المعادلة التالية بطريقة RK4
			\[
			p' = 1.1\,p, \quad p(0)=2
			\]
			على الفترة \([0,5]\) مع التقسيم:
			\[
			t_0=0,\; t_1=1,\; t_2=2,\; t_3=3,\; t_4=4,\; t_5=5
			\]
			وبالتالي فإن:
			\[
			h = 1
			\]
			نعرّف:
			\[
			f(t,p) = 1.1\,p
			\]
			\textbf{الخطوة الأولى: حساب \( p_1 \)}
			\[
			k_1 = f(0,2) = 2.2
			\]
		\end{frame}
		
		\begin{frame}
			
			\[
			k_2 = f\left(0.5,\, 2 + 0.5 \times 2.2 \right) = f(0.5, 3.1) = 3.41
			\]
			\[
			k_3 = f\left(0.5,\, 2 + 0.5 \times 3.41 \right) = f(0.5, 3.705) = 4.0755
			\]
			\[
			k_4 = f(1,\, 2 + 4.0755) = f(1, 6.0755) = 6.68305
			\]
			\[
			p_1 = 2 + \frac{1}{6}(2.2 + 2(3.41) + 2(4.0755) + 6.68305)
			\]
			\[
			p_1 \approx 2 + \frac{1}{6}(23.85305) \approx 5.9755
			\]
			\textbf{الخطوة الثانية: حساب \( p_2 \)}
			\[
			k_1 = 1.1(5.9755) = 6.57305
			\]
			\[
			k_2 = 1.1(5.9755 + 0.5 \times 6.57305) = 10.18823
			\]
			\[
			k_3 = 1.1(5.9755 + 0.5 \times 10.18823) = 12.16723
			\]
		\end{frame}
		
		\begin{frame}
					\[
			k_4 = 1.1(5.9755 + 12.16723) = 19.9560
			\]
			\[
			p_2 = 5.9755 + \frac{1}{6}(6.57305 + 2(10.18823) + 2(12.16723) + 19.9560)
			\]
			\[
			p_2 \approx 17.854
			\]
			وبنفس الطريقة نجد بقية القيم
		\end{frame}
		
		\begin{frame}{مقارنة النتائج}
			\begin{table}[h]
				\addcontentsline{toc}{section}{مقارنة النتائج}
				\centering
				\renewcommand{\arraystretch}{1.5}
				\begin{english}
					\centering
					\begin{tabular}{|c|c|c|c|}
					\hline
					$t$ & $p_{\text{DTM}}(t)$ & $p_{\text{RK4}}(t)$ & Exact \\ 
					\hline
					0 & 2.00000000 & 2.00000000 & 2.00000000 \\ 
					1 & 6.00833205 & 5.97567500 & 6.00833205 \\ 
					2 & 18.05002700 & 17.85434585 & 18.05002700 \\ 
					3 & 54.22527782 & 53.34588407 & 54.22527784 \\ 
					4 & 162.90173570 & 159.38883290 & 162.90173730 \\ 
					5 & 489.38368110 & 476.22793210 & 489.38386460 \\ 
					\hline
				\end{tabular}
				\end{english}
				\caption{\centering مقارنة بين طريقة التحويل التفاضلي (DTM) وطريقة رونج-كوتا من الرتبة الرابعة (RK4) مع الحل الدقيق للمعادلة التفاضلية $p'=1.1p$}
			\end{table}
		\end{frame}
		
		\begin{frame}{الاستنتاج}\doublespacing
			
			في نهاية البحث نستنتج ان طريقةالتحويل التفاضلي طريقة جيدة جداً لإيجاد حل دالي تقريبي لحل المعادلة التفاضلية حيث يتم ايجاد معاملات المتسلسلة من خلال علاقة تكرارية تكون بسيطة حسابياً ، ولكن لكي ان نصل الى حل أدق يجب ان نزيد في عدد الحدود المحسوبة.
		\end{frame}
		
		\begin{frame}
			\Huge\centering
			\textbf{شكراً لحسن استماعكم}
		\end{frame}
\end{document}